\section{相关化合物的合成方法}
\label{sec:neighbor-routes}

本节按成相过程比较相关化合物，不按论文逐篇罗列。图~\ref{fig:route-matrix}概括可比较的实验项目，表~\ref{tab:synthesis-conditions}保留数值和文献出处。以下只讨论相关研究提供的方法信息，不把任何路线改写成Ba$_5$Y$_{12}$Zn[O(SiO$_4$)]$_8$的候选配方。

\subsection{高温溶液路线}

Y--Si--O文献首先说明，“长出晶体”不等于“结构正确”。Na$_2$MoO$_4$助熔得到的$\beta$-Y$_2$Si$_2$O$_7$伴有其他硅酸盐晶体，最终相型由单晶结构测定确认；这类研究可参照营养物与助熔剂、坩埚盖合、慢冷和分离，但不能给出目标相的稳定范围\cite{redhammer2003beta}。$^{89}$Y MAS--NMR重访Y$_2$Si$_2$O$_7$/Y$_2$SiO$_5$多型表明，局域探针可补充衍射区分多型，却不能单独证明目标相成相\cite{becerro2004revisiting}。Y--Si--O合成综述和经典稀土硅酸盐晶体化学分类可统一多型和硅酸根聚合度术语，但不能替代原始实验记录\cite{sun2014recent,felsche1973the}。

在Ba--Zn--Si--O体系中，开放体系高温溶液研究只确认Ba$_3$Zn$_4$Si$_4$O$_{15}$的结构和路线类别，精确批料与热史仍记为\texttt{NA}\cite{ababaikeri2024ba3zn4si4o15}。Zn$_2$SiO$_4$的历史熔融溶液研究表明，慢冷、晶体分离和助熔成分必须同时记录；Pb/F条件还涉及毒性、挥发和废物处置审查，不能作为默认选择\cite{giess1982zn2sio4}。来源明确的无坩埚熔体和Y$_2$SiO$_5$/Y$_2$Si$_2$O$_7$高温结晶研究，可参照容器接触、自发成核和结构鉴定，但不能证明目标相具有相同的熔体生长性\cite{christensen1994investigation,leonyuk1999high}。另一条Leonyuk记录只有题名证据，只能确认单晶对象和结构精修，不能据此判断合成方法\cite{leonyuk1999crystal}。

\textbf{对目标实验的启示。}应将投料和助熔剂分开记录，注明容器接触，并连续记录冷却和分离过程。晶体结构、批量衍射和成分分析需要配套进行。相关化合物的具体数值、Pb/F体系的适用性以及由晶体外形推定目标相身份，均不能直接外推。

\subsection{固相成相路线}

BaMSiO$_4$（M=Co、Zn、Mg）经单晶X射线和粉末中子衍射建立stuffed-tridymite同构族，说明同价四面体位替代可以保持母拓扑，但该骨架不能等同于目标相的零维混合构筑\cite{liu1993structures}。单斜Ba$_2$ZnSi$_2$O$_7$的单晶精修和单斜Ba$_2$MgSi$_2$O$_7$的结构报告分别提供Zn、Mg参照；它们可用于识别竞争相和检查替代风险，但不能据此声称与目标相同型\cite{kaiser2002crystal,aitasalo2006crystal}。

受控取代系列可示范如何成组记录组成、结构、热史和表征。Ba/Sr取代物由SCXRD与EDX联合确认；BaZn$_{2-x}$Co$_x$Si$_2$O$_7$及Ba/Sr--Zn--Si/Ge系列的HT-XRD显示组成相关的结构保持区和相边界。这些研究说明，名义替代量必须与实际组成和原位结构结果对应，不能跨体系量化目标相的Y/Zn容限\cite{thieme2015ba1,kerstan2013bazn2si2o7,thieme2016negative}。BaZn$_2$Si$_2$O$_7$固溶体综述可整理热膨胀随组成变化的背景，但不替代单个原始烧结批次\cite{thieme2022solid}。含Ba$_{1-x}$Sr$_x$Zn$_2$Si$_2$O$_7$的玻璃陶瓷研究同样只提供功能和晶化语境，不能作为目标相的直接合成依据\cite{thieme2017variable}。多气氛Ba$_2$ZnSi$_2$O$_7$陶瓷对照表明，气氛变化必须结合Zn损失、副相和表面价态解释，不能只凭介电性能变化判断结构稳定\cite{zou2019anti}。

\textbf{对目标实验的启示。}应并列记录化学计量称量、复磨复烧、气氛和实际组成，并采用PXRD/Rietveld、成分分析及必要的HT-XRD。相关研究中的同价替代、热膨胀或介电性能结果，不能直接写成目标相的异价替代机制、相纯范围或还原稳定性。

\subsection{活化与提拉}

机械化学研究只支持“预活化可能改变后续成相路径”这一认识。钇氧磷灰石研究给出空气中256\,rpm、1--3\,h以及玛瑙罐和玛瑙球的机械处理条件；处理后仍需热处理，且多个Y--Si--O相共存。这些数值只描述该混相体系，不能转写为目标相的活化条件\cite{tzvetkov2001effects}。Y$_2$SiO$_5$粉体论文只明确报道LiYO$_2$辅助的固--液相烧结以及显微、XRD和SEM--EDS表征，不能因对象是粉体就称为机械化学\cite{yldrm2009y2sio5}。

Czochralski研究提供熔体、Ir坩埚、气氛、SiO$_2$挥发、提拉与旋转以及退火等参照。Y$_2$SiO$_5$的生长、形貌和光谱联测说明，晶体生长后仍需进行结构和成分鉴定。稀土正硅酸盐的历史提拉研究及后续学位论文可用于规范工艺记录，但目前没有目标相的提拉生长证据\cite{pang2005study,shoudu1999czochralski,brandle1986czochralski,alizadeh2021spectroscopy}。

\textbf{对目标实验的启示。}可参照机械活化过程中的污染检查，以及提拉法对熔体、坩埚和挥发的控制；同时必须验证批量相组成、实际元素比、夹杂和单晶结构。普通混匀不能称为机械化学，Y$_2$SiO$_5$/Ln$_2$SiO$_5$的可提拉性也不等同于目标相的可提拉性。所有数值均以表~\ref{tab:synthesis-conditions}的原始文献记录为限。
